\chapter{功能、可靠性与性能验证}
\label{chap:validation}

\section{证据分层}

我们不用单一“tests passed”概括整个系统。不同层的证据回答不同问题，下层通过不会自动升格为上层结论。

\begin{table}[htbp]
  \centering
  \caption{AgentOS-uCore 的六层验证证据}
  \label{tab:validation-layers}
  \begin{tabularx}{\textwidth}{L{2.7cm}Y Y}
    \toprule
    证据层 & 能够回答 & 不能回答 \\
    \midrule
    Focused checker & 关键调用顺序、UAPI 布局、生产对象边界是否漂移 & 真实 RISC-V syscall 语义 \\
    Host/mutation 模型 & 账户、ring、resync、relay schema 与错误输入是否 fail closed & Guest 等待、调度和 VFS 实际行为 \\
    Cross build/link & 当前生产对象能否在 RV64 编译链上链接 & 程序运行时会走到正确路径 \\
    Stack/module gate & 调用图的栈预算和生产模块依赖是否越界 & 全部时序竞态已被排除 \\
    QEMU Guest & identity、Context、V2/V3、IPC、VFS、wait/wake 在实际 Guest 中是否闭环 & 其他硬件或任意工作负载的泛化 \\
    配对性能 & 指定版本、工具链、QEMU 和 Host 上的配对工程表现 & 某个内核机制的单独因果效应 \\
    \bottomrule
  \end{tabularx}
\end{table}

\section{功能验证矩阵}

\begin{longtable}{L{1.2cm}L{3.0cm}L{5.6cm}L{3.1cm}}
  \caption{六项任务的关键功能与证据入口}\label{tab:function-matrix}\\
  \toprule
  任务 & 核心约束 & 实现证据 & 动态验收焦点 \\
  \midrule
  \endfirsthead
  \toprule
  任务 & 核心约束 & 实现证据 & 动态验收焦点 \\
  \midrule
  \endhead
  一 & 内核可识别 Agent，不信任用户自报 & PCB identity、trusted exec、Context mirror、lifecycle 与资源账户 & create、fork、exec、exit、generation reuse、关闭门禁 \\
  二 & 结构化参数先验证后副作用 & 25 工具目录、typed V2、compact batch、V3 contract、Task SQ/CQ & 错误 schema/capability/provenance/deadline 拒绝；96 个 sequence 性能表 \\
  三 & Context 和 Evidence 不得被用户直改 & 7-page Context、rollback、provenance、Evidence Ring、workflow fence & push、query、rollback、critical 保留、seal 幂等 \\
  四 & 查询走真实 VFS 和当前 boot 索引 & 显式 metadata、scan/index 一致性、typed watch、resync generation & create、unlink、incarnation、scan/index 结果相同、3$\times$4$\times$15 负载 \\
  五 & 等待不忙轮询，线程数不放大 workflow 份额 & event queue、LLM correlation、heartbeat、workflow EEVDF & timeout/cancel、ready-age、504 exact wake、Jain fairness \\
  六 & 同一 AgentOS Guest 的 traversal/indexed 工作量对齐，产品闭环与性能基准分工 & deterministic \code{labdemo_ucore}、Nexus、Host validators、AB/BA runner & 16 paired boot；Nexus 3 turn/4 identity/failure/replan/deny \\
  \bottomrule
\end{longtable}

\section{可靠性与负路径}

\subsection{静态合同与 mutation}

Focused checker 不只搜索“函数名存在”，还约束生产调用链和禁止的旁路。例如，workflow fence checker 锁定唯一 syscall cut 和 metadata/live-query、VFS、credit、evidence 的顺序；Nexus checker 锁定 System no-input 路径不读 TASK capsule、Research/Analyst 验证 handle 和 provenance、副作用审批的全字段绑定。mutation tests 改动长度、版本、顺序、digest、lifecycle 或身份字段，检查系统是否拒绝。

\subsection{构建与栈安全}

\code{agent-uapi-check} 对齐内核和用户头文件的 size/offset；\code{agent-module-check} 检查真实生产 object list 和依赖；\code{kernel-stack-check} 与 \code{user-stack-check} 在实际编译调用图上检查栈深。Nexus Guest 使用大型 static buffer 和专用小结构 history，避免将 24 个 TASK event 嵌入每个 turn 而突破 uCore 当前 \code{MAXFILE=274432} bytes。Image target 在 mkfs 前以 raw binary 大小 fail closed，不通过扩大 FS ABI 规避问题。

\subsection{QEMU 负路径}

下列结果是产品语义的一部分，不是可以从报告中删除的“异常噪声”：

\begin{itemize}
  \item malformed typed KV 在工具运行前返回 \code{BAD_PARAM}；
  \item capability、route、scope 或 provenance 不符合时返回 \code{DENIED}；
  \item 旧 generation handle 或已消费 response 返回 \code{STALE}；
  \item event/LLM/task deadline 超时返回 \code{TIMEOUT}，而不无限忙轮询；
  \item artifact 容量或 MESSAGE source queue 受限时返回 \code{NO_SPACE}；
  \item 审批 deny 作为结构化 tool result 回灌，不终止整个 Nexus session；
  \item observer 正常关闭前经过 final drain/join/EOF，关闭之后的任何 frame 都使 strict validator 失败。
\end{itemize}

\section{Nexus 严格 replay 验收}

\subsection{方法}

\code{agentos-nexus-replay} 不是 Host 重放一段录屏，而是启动当前 RV64 镜像、连接 controller/observer，将每条 fixture response 与 Guest 当次产生的 canonical \code{MODEL_REQUEST} SHA-256 绑定。请求文本、history 或稳定 tool projection 发生任何变化，原 fixture 都会 fail closed。

Validator 同时检查 ready-before-data 与 active request 不复用。因果顺序必须为 \code{request} $<$ \code{response} $<$ effects $<$ next request。每个 TASK 只有一个 terminal，parent task 必须在当次 transcript 中存在；controller 与 observer 的 TASK 安全 metadata 子序列一致，kernel audit enqueue/consume 成对，Context sequence 不与 audit sequence 混淆。

\subsection{黄金路径}

三个连续用户 turn 覆盖：

\begin{enumerate}
  \item System 执行无输入本 boot 快照，生成真实 System artifact；
  \item Research 先因 stale/missing handle 结构化失败，Coordinator 创建新 task 重新规划，后续从本地 measurement capsule 生成 Research artifact；
  \item Analyst 读取两份已验 artifact，生成带 handle/digest/source 的 report；发布请求被人工拒绝，但最终答案仍安全完成。
\end{enumerate}

该路径要求 worker 在等待时真实进入 \code{WAITING_EVENT}，Coordinator 等待 provider 时进入 \code{WAITING_LLM}。成功路径不允许 Host 选择工具、读 Guest 业务文件或伪造 artifact。严格 replay 能复验此闭环；它不等于 DeepSeek live provider 已在该次测试中成功。

\section{一次性性能数据集}

\subsection{数据来源与独立单位}

规范运行位于 \code{one_shot_metrics/data/20260811}，由 \code{CANONICAL_RUN} 指向。采集共使用 30 个 fresh QEMU boot，保留 33 个 raw 文件，生成 19 张长表和 7,498 行结构化记录。每行保留 raw 相对路径、来源 SHA-256 和来源内序号，因而图表可回到原始串口证据。

\begin{table}[htbp]
  \centering
  \caption{一次性采集的实验族与正确独立单位}
  \label{tab:one-shot-families}
  \begin{tabularx}{\textwidth}{L{3.0cm}L{4.0cm}L{3.2cm}Y}
    \toprule
    数据族 & 设计 & 独立单位 & 关键限定 \\
    \midrule
    traversal/indexed & 16 个平衡 AB/BA paired boot & boot & core 与 E2E 分开报告 \\
    file catalog grid & $3\times4\times15=180$ inner pair & boot 为推断单位 & 同一 hit-count 只有 1 个 source boot \\
    compact batch/scalar V3/SQ-CQ & 4 boot $\times$ 8 rotating round & boot/block & 每序列 16 op，共 96 sequence \\
    workflow EEVDF & 6 boot，并发 1--4 & boot/window 或 wake probe & 504 exact probe，不从桶插值 \\
    \bottomrule
  \end{tabularx}
\end{table}

同一 boot 中的 inner pair 是重复测量，不是独立机器实验。本次数据定位为工程探索与文档绘图，不进行跨硬件总体推断。

\section{Core 与端到端窗口}

16 个独立 paired boot 中，indexed core 在 16/16 配对中更快；traversal/indexed core 中位数分别为 34.713 ms 和 13.294 ms，配对中位加速比为 3.118$\times$。但端到端窗口仅 3/16 个配对 indexed 更快，\code{indexed - traversal} 的中位数为 $+13.452$ ms。因此我们只宣称 core path 在该负载下获益，不宣称整个系统窗口已加速。

AB/BA 顺序分层的 core speedup 中位数为 1.4535$\times$ 和 5.4805$\times$。差异提示顺序、缓存或其他 boot 层因素仍然很强，正是我们保留平衡配对和原始分层的原因。

\section{File query catalog 工作量}

目录规模 24/64/96 与 exact hit count 1/2/4/8 形成 12 格，每格 15 个 AB/BA inner pair。图\ref{fig:catalog-speedup}与图\ref{fig:scan-groups}展示索引对不同选择性的影响。图中 first scan 是显式 warmup 后的第一次 timed scan，不被伪称为可证明的物理冷缓存。同一 hit-count 只有一个 source boot，所以热力图是工程地形，不提供 15 个独立 boot 的置信结论。

\section{Task 通道序列分布}

\ReportFigure[0.95\textwidth]
  {../figures/performance/04_task_latency_distributions.pdf}
  {../figures/performance/04_task_latency_distributions.pdf}
  {4 个 boot、8 个旋转顺序下 96 个 16-op sequence 延迟分布}
  {fig:task-latency}

Compact batch、scalar V3 与 SQ/CQ 的 16-op sequence 中位数分别为 561 us、2051 us 和 1620.5 us。这是完整 sequence 的 Host/Guest 测量，不应将 10 ms tick 粒度的 \code{service_start_interval_tick} 均摊成伪造的单请求微秒延迟。此次数据显示了不同路径的产品代价，而不是声称异步通道必然快于同步 syscall。

\section{EEVDF 唤醒与公平性}

图\ref{fig:eevdf-performance}同时展示 exact wake ECDF 与 Jain fairness。在 504 个带完整前缀的 probe 中，425 个 ready age 为 0 tick，79 个为 1 tick。EEVDF boot 05 有 3 条并发串口 sample summary 拼接；提取器保留 warning 并排除不完整 summary，504 条 exact probe 全部保留，没有插值或静默修补。

并发 1 的 Jain=1 是平凡基线；scenario 16 的四波不是 16-way physical concurrency。我们在报告中保留这些限制，防止图形的美观性取代实验定义。

\section{结果解释的最小规则}

\begin{evidencebox}[引用性能数字时必须同时披露]
对外引用必须给出 target、source commit、工具链、QEMU/Host 环境、样本数、单位、聚合方式和失败/warning。paired boot 才是 core/E2E 比较的独立单位；inner pair 不能冒充 fresh boot。Nexus 的 \code{nexus_meas} 只是历史快照，它不能替代新的 \code{contest-demo} 采集。
\end{evidencebox}
